% Options for packages loaded elsewhere
\PassOptionsToPackage{unicode}{hyperref}
\PassOptionsToPackage{hyphens}{url}
\documentclass[
]{article}
\usepackage{xcolor}
\usepackage{amsmath,amssymb}
\setcounter{secnumdepth}{5}
\usepackage{iftex}
\ifPDFTeX
  \usepackage[T1]{fontenc}
  \usepackage[utf8]{inputenc}
  \usepackage{textcomp} % provide euro and other symbols
\else % if luatex or xetex
\fi
\usepackage{lmodern}
\ifPDFTeX\else
  % xetex/luatex font selection
\fi
% Use upquote if available, for straight quotes in verbatim environments
\IfFileExists{microtype.sty}{% use microtype if available
  \usepackage[]{microtype}
  \UseMicrotypeSet[protrusion]{basicmath} % disable protrusion for tt fonts
}{}
\makeatletter
\@ifundefined{KOMAClassName}{% if non-KOMA class
  \IfFileExists{parskip.sty}{%
    \usepackage{parskip}
  }{% else
    \setlength{\parindent}{0pt}
    \setlength{\parskip}{6pt plus 2pt minus 1pt}}
}{% if KOMA class
  \KOMAoptions{parskip=half}}
\makeatother
\usepackage{listings}
\newcommand{\passthrough}[1]{#1}
\lstset{defaultdialect=[5.3]Lua}
\lstset{defaultdialect=[x86masm]Assembler}
\usepackage{longtable,booktabs,array}
\newcounter{none} % for unnumbered tables
\usepackage{calc} % for calculating minipage widths
% Correct order of tables after \paragraph or \subparagraph
\usepackage{etoolbox}
\makeatletter
\patchcmd\longtable{\par}{\if@noskipsec\mbox{}\fi\par}{}{}
\makeatother
\ifLuaTeX
\else
\fi
\ifLuaTeX
\fi
\setlength{\emergencystretch}{3em} % prevent overfull lines
\providecommand{\tightlist}{%
  \setlength{\itemsep}{0pt}\setlength{\parskip}{0pt}}
% ═══ 由 环境/pdf/latex/md到LaTeX.py 生成 ═══
% 中文字体：思源宋体（子集，SIL OFL 1.1）；拉丁/数学：Latin Modern（随 TeX Live 分发）
% 编译：xelatex（本仓库自带 wasm 版 XeTeX 引擎，见 环境/pdf/latex/编译LaTeX.mjs）
\usepackage[T1]{fontenc}      % 否则 ASCII 的 < > | 在 OT1 下排成 ¡ ¿ —
\usepackage{lmodern}          % 必需：数学的物理字模（否则 xdvipdfmx 拒绝出 PDF）
\usepackage{amsmath,amssymb,mathtools}
\usepackage{booktabs,longtable,array,multirow}
\usepackage{graphicx}
\usepackage{xcolor}
\usepackage{listings}
% 不用 hyperref：新版 hyperref 强制依赖 bookmark.sty，而 bundle 里没有。
% 链接命令用降级定义，保证正文里的 \href / \url 不会报错。
\providecommand{\href}[2]{#2}
\providecommand{\url}[1]{\texttt{#1}}
\providecommand{\hypertarget}[2]{#2}
\providecommand{\hyperlink}[2]{#2}
\providecommand{\urlstyle}[1]{}
\providecommand{\texorpdfstring}[2]{#1}   % hyperref 的命令，剥离后仍被模板/标题用到
\providecommand{\phantomsection}{}
\providecommand{\pdfstringdef}[2]{}
\providecommand{\hypersetup}[1]{}

% ── 三套字体：中文 / 符号 / emoji（按字符 cmap 自动选择）──
\font\zhfont="[NotoSerifSC-Regular.ttf]:script=hani" at 10.5pt
\font\symfont="[DejaVuSans.ttf]" at 10.5pt
\font\emofont="[NotoEmoji-Regular.ttf]" at 10.5pt
\font\zhmonofont="[NotoSansSC-Regular.ttf]" at 9pt
% 用法：\zh{中文} / \sym{→} / \emo{🦙}
% 注意：必须**带参数**。写成 \def\zh{{\zhfont}} 是错的——字体赋值只在小群里有效，
% 那个小群随 \zh 立刻结束，于是中文仍然用拉丁字体排，报 Missing character。
\def\zh#1{{\zhfont #1}}
\def\sym#1{{\symfont #1}}
\def\emo#1{{\emofont #1}}
% 含汉字的行内代码：listings 的 \lstinline 遇到汉字必炸（`\lst@arg ->判`
% 未定义控制字，catcode 设 12 也无效），所以直接自己排。
\def\codezh#1{{\zhmonofont #1}}

\def\zhglue{\hskip0pt plus .06em minus .01em}
\linespread{1.12}
\setlength{\parindent}{0pt}
\setlength{\parskip}{0.45em}

% ── 代码块：等宽 + 中文字体（listings 默认字体不含汉字）──
\lstset{
  basicstyle=\zhmonofont,
  breaklines=true,
  breakatwhitespace=false,
  columns=fullflexible,
  keepspaces=true,
  showstringspaces=false,
  frame=single,
  framesep=4pt,
  rulecolor=\color{gray!45},
  backgroundcolor=\color{gray!7},
  xleftmargin=6pt,
  extendedchars=true,
  tabsize=2,
  % 代码块里的 emoji：NotoSansSC 没有这些字形（listings 不做字符级回退）。
  % 用 escapeinside 开口子，转换器把 emoji 逐字包成 (*@{\emo{⭐}}@*)。
  % 试过 literate={⭐}{{\emofont ⭐}}1，会把 ASCII 字母也吞掉并报
  % `Improper alphabetic constant`，不可用。
  escapeinside={(*@}{@*)},
}
\renewcommand{\lstlistingname}{\zh{代\zhglue 码}}

% ── 表格线宽（booktabs 风格）──
\renewcommand{\arraystretch}{1.25}

% ── 标题样式 ──
\usepackage{titlesec}
\titleformat{\section}{\Large\bfseries}{\thesection}{0.6em}{}
\titlespacing*{\section}{0pt}{1.1em}{0.5em}
\urlstyle{same}
\hypersetup{
  pdftitle={\zh{课\zhglue 题}10-\zh{长\zhglue 程}Agent},
  pdflang={zh-CN},
  hidelinks,
  pdfcreator={LaTeX via pandoc}}

\title{\zh{课\zhglue 题}10-\zh{长\zhglue 程}Agent}
\author{}
\date{}
\begin{document}
\maketitle

{
\setcounter{tocdepth}{2}
\tableofcontents
}
\section{\zh{课\zhglue 题}10 \zh{开\zhglue 题\zhglue 简\zhglue 报} \sym{·} \zh{长\zhglue 程} Agent \zh{与} self-judge}\label{ux8bfeux989810-ux5f00ux9898ux7b80ux62a5-ux957fux7a0b-agent-ux4e0e-self-judge}

\begin{quote}
\zh{模\zhglue 块\zhglue ：}\textbf{M8} \zh{｜} \zh{周\zhglue 次\zhglue ：}\textbf{W14\zh{（}12-14\sym{→}12-20\zh{）}} \zh{｜} \zh{算\zhglue 力\zhglue ：}\textbf{T2 Kaggle 2\sym{×}T4\zh{，\zhglue 计\zhglue 划} 8 GPU\sym{·}\zh{小\zhglue 时}} \zh{唐\zhglue 杰\zhglue 作\zhglue 业\zhglue 对\zhglue 应\zhglue ：}\textbf{\sym{⑤} \zh{训\zhglue 长\zhglue 程} Agent\zh{，\zhglue 鼓\zhglue 励} self-judge loop} \zh{｜} \zh{判\zhglue 分\zhglue 来\zhglue 源\zhglue ：}\textbf{\zh{自\zhglue 建}}\zh{（\zhglue 消\zhglue 融\zhglue 实\zhglue 验} + \zh{轨\zhglue 迹\zhglue 日\zhglue 志\zhglue ）} \zh{手\zhglue 写\zhglue ：}\textbf{R10\zh{（}self-judge \zh{协\zhglue 议\zhglue 设\zhglue 计\zhglue ）}}
\end{quote}

\subsection{\zh{一\zhglue 、\zhglue 研\zhglue 究\zhglue 背\zhglue 景}}\label{ux4e00ux7814ux7a76ux80ccux666f}

\zh{从}''\zh{回\zhglue 答\zhglue 问\zhglue 题}''\zh{到}''\zh{完\zhglue 成\zhglue 任\zhglue 务}''\zh{的\zhglue 跨\zhglue 越\zhglue ，\zhglue 靠\zhglue 的\zhglue 是} \textbf{harness\zh{（\zhglue 外\zhglue 壳\zhglue ）}}\zh{：\zhglue 工\zhglue 具\zhglue 调\zhglue 用\zhglue 、\zhglue 上\zhglue 下\zhglue 文\zhglue 管\zhglue 理\zhglue 、\zhglue 失\zhglue 败\zhglue 重\zhglue 试\zhglue 、\zhglue 自\zhglue 我\zhglue 评\zhglue 估\zhglue 。} Karpathy \zh{与\zhglue 李\zhglue 宏\zhglue 毅\zhglue 都\zhglue 专\zhglue 讲} \textbf{context engineering}\zh{；}nanochat \zh{里} \passthrough{\lstinline!execution.py!} \zh{就\zhglue 是\zhglue 一\zhglue 个\zhglue 最\zhglue 小\zhglue 的}''\zh{能\zhglue 执\zhglue 行\zhglue 代\zhglue 码\zhglue 的} Agent \zh{外\zhglue 壳}''\zh{。} \zh{本\zhglue 课\zhglue 题\zhglue 的\zhglue 落\zhglue 点\zhglue 是}\textbf{\zh{长\zhglue 程\zhglue 任\zhglue 务}}\zh{：\zhglue 多\zhglue 步\zhglue 串\zhglue 联\zhglue 后\zhglue 误\zhglue 差\zhglue 累\zhglue 积\zhglue ，\zhglue 以\zhglue 及}\textbf{self-judge\zh{（\zhglue 自\zhglue 我\zhglue 评\zhglue 判\zhglue ）\zhglue 到\zhglue 底\zhglue 有\zhglue 没\zhglue 有\zhglue 用}}\zh{。}

\subsection{\zh{二\zhglue 、\zhglue 研\zhglue 究\zhglue 问\zhglue 题}}\label{ux4e8cux7814ux7a76ux95eeux9898}

\begin{quote}
\textbf{\zh{多\zhglue 步\zhglue 任\zhglue 务\zhglue 中\zhglue ，\zhglue 失\zhglue 败\zhglue 主\zhglue 要\zhglue 来\zhglue 自\zhglue 模\zhglue 型\zhglue 能\zhglue 力\zhglue 还\zhglue 是\zhglue 来\zhglue 自} harness \zh{设\zhglue 计\zhglue ？\zhglue 加\zhglue 上} self-judge loop \zh{后\zhglue 成\zhglue 功\zhglue 率\zhglue 提\zhglue 升\zhglue 多\zhglue 少}------\zh{还\zhglue 是\zhglue 反\zhglue 而\zhglue 下\zhglue 降\zhglue ？}}
\end{quote}

\subsection{\zh{三\zhglue 、\zhglue 攻\zhglue 关\zhglue 线\zhglue 索}}\label{ux4e09ux653bux5173ux7ebfux7d22}

\begin{enumerate}
\def\labelenumi{\arabic{enumi}.}
\tightlist
\item
  \textbf{\zh{最\zhglue 小} ReAct}\zh{：\zhglue 不\zhglue 用\zhglue 任\zhglue 何\zhglue 框\zhglue 架\zhglue ，\zhglue 手\zhglue 写} \passthrough{\codezh{Thought {\sym{→}} Action {\sym{→}} Observation}} \zh{循\zhglue 环\zhglue （\zhglue 三\zhglue 四\zhglue 十\zhglue 行\zhglue 就\zhglue 能\zhglue 跑\zhglue ）}
\item
  \textbf{\zh{工\zhglue 具\zhglue 设\zhglue 计}}\zh{：\zhglue 先\zhglue 给\zhglue 两\zhglue 个\zhglue 工\zhglue 具\zhglue （\zhglue 计\zhglue 算\zhglue 器} / \zh{受\zhglue 限\zhglue 代\zhglue 码\zhglue 执\zhglue 行\zhglue ）\zhglue 。}\textbf{\zh{工\zhglue 具\zhglue 接\zhglue 口\zhglue 的\zhglue 清\zhglue 晰\zhglue 度\zhglue 往\zhglue 往\zhglue 比\zhglue 模\zhglue 型\zhglue 智\zhglue 力\zhglue 更\zhglue 决\zhglue 定\zhglue 成\zhglue 败}}
\item
  \textbf{\zh{长\zhglue 程\zhglue 任\zhglue 务\zhglue 构\zhglue 造}}\zh{：\zhglue 把\zhglue 课\zhglue 题}09 \zh{的\zhglue 题\zhglue 改\zhglue 造\zhglue 成}''\zh{需\zhglue 要} 3--5 \zh{步}''\zh{的\zhglue 复\zhglue 合\zhglue 任\zhglue 务\zhglue （\zhglue 如\zhglue ：\zhglue 先\zhglue 算} \sym{→} \zh{再\zhglue 验} \sym{→} \zh{再\zhglue 格\zhglue 式\zhglue 化\zhglue ）}
\item
  \textbf{\zh{失\zhglue 败\zhglue 分\zhglue 类}}\zh{：\zhglue 格\zhglue 式\zhglue 错} / \zh{工\zhglue 具\zhglue 调\zhglue 错} / \zh{目\zhglue 标\zhglue 漂\zhglue 移\zhglue （\zhglue 忘\zhglue 了\zhglue 原\zhglue 始\zhglue 要\zhglue 求\zhglue ）}/ \zh{循\zhglue 环\zhglue 打\zhglue 转} / \zh{提\zhglue 前\zhglue 放\zhglue 弃}
\item
  \textbf{self-judge \zh{的\zhglue 两\zhglue 种\zhglue 形\zhglue 态}}\zh{：}

  \begin{itemize}
  \tightlist
  \item
    \textbf{\zh{自\zhglue 评\zhglue 重\zhglue 试}}\zh{：\zhglue 模\zhglue 型\zhglue 给\zhglue 自\zhglue 己\zhglue 打\zhglue 分\zhglue ，\zhglue 低\zhglue 于\zhglue 阈\zhglue 值\zhglue 就\zhglue 重\zhglue 做\zhglue （}rasbt ch5 \zh{的} self-refinement\zh{）}
  \item
    \textbf{\zh{自\zhglue 评\zhglue 过\zhglue 滤}}\zh{：\zhglue 生\zhglue 成\zhglue 多\zhglue 个\zhglue 候\zhglue 选\zhglue ，\zhglue 自\zhglue 己\zhglue 挑\zhglue 最\zhglue 好\zhglue 的\zhglue （}best-of-N \zh{的\zhglue 变\zhglue 体\zhglue ）}
  \end{itemize}
\item
  \textbf{\zh{关\zhglue 键\zhglue 实\zhglue 验\zhglue （\zhglue 消\zhglue 融\zhglue ）}}\zh{：\zhglue 同\zhglue 样\zhglue 的\zhglue 任\zhglue 务\zhglue 集\zhglue ，\zhglue 三\zhglue 组\zhglue 对\zhglue 照} = \zh{无} self-judge / \zh{自\zhglue 评\zhglue 重\zhglue 试} / \zh{自\zhglue 评\zhglue 过\zhglue 滤}
\item
  \textbf{\zh{警\zhglue 惕}}\zh{：}self-judge \textbf{\zh{会\zhglue 放\zhglue 大\zhglue 原\zhglue 有\zhglue 错\zhglue 误}}\zh{（\zhglue 模\zhglue 型\zhglue 判\zhglue 不\zhglue 出\zhglue 自\zhglue 己\zhglue 错\zhglue 在\zhglue 哪\zhglue ，\zhglue 反\zhglue 而\zhglue 确\zhglue 信\zhglue ）}\sym{→} \zh{必\zhglue 须\zhglue 留\zhglue 失\zhglue 败\zhglue 案\zhglue 例}
\end{enumerate}

\subsection{\zh{四\zhglue 、\zhglue 里\zhglue 程\zhglue 碑}}\label{ux56dbux91ccux7a0bux7891}

\begin{itemize}
\tightlist
\item[$\square$]
  M1 \zh{手\zhglue 写\zhglue 最\zhglue 小} ReAct \zh{循\zhglue 环\zhglue （\zhglue 无\zhglue 框\zhglue 架\zhglue ）\zhglue ，\zhglue 跑\zhglue 通\zhglue 单\zhglue 步\zhglue 任\zhglue 务}
\item[$\square$]
  M2 \zh{加\zhglue 入\zhglue 第\zhglue 二\zhglue 个\zhglue 工\zhglue 具} + \zh{多\zhglue 步\zhglue 任\zhglue 务\zhglue 集\zhglue （}\sym{≥}30 \zh{个\zhglue 任\zhglue 务\zhglue ）}
\item[$\square$]
  M3 \textbf{\zh{轨\zhglue 迹\zhglue 记\zhglue 录\zhglue 规\zhglue 范}}\zh{：\zhglue 每\zhglue 步\zhglue 的} thought/action/observation \zh{落\zhglue 盘\zhglue （\zhglue 可\zhglue 复\zhglue 现\zhglue ）}
\item[$\square$]
  M4 \zh{基\zhglue 线\zhglue 成\zhglue 功\zhglue 率} + \textbf{\zh{失\zhglue 败\zhglue 五\zhglue 分\zhglue 类}}\zh{（\zhglue 每\zhglue 类\zhglue 给} 2 \zh{个\zhglue 真\zhglue 实\zhglue 案\zhglue 例\zhglue ）}
\item[$\square$]
  M5 \textbf{self-judge \zh{消\zhglue 融}}\zh{：\zhglue 三\zhglue 组\zhglue 对\zhglue 照\zhglue ，\zhglue 给\zhglue 出\zhglue 成\zhglue 功\zhglue 率\zhglue 差\zhglue 与\zhglue 误\zhglue 差\zhglue 棒}
\item[$\square$]
  M6 \textbf{\zh{反\zhglue 例\zhglue 收\zhglue 集}}\zh{：\zhglue 找\zhglue 出} \sym{≥}3 \zh{个}''self-judge \zh{导\zhglue 致\zhglue 更\zhglue 差}''\zh{的\zhglue 案\zhglue 例\zhglue 并\zhglue 解\zhglue 释\zhglue 机\zhglue 制}
\item[$\square$]
  M7 \zh{一\zhglue 页\zhglue 报\zhglue 告} + \zh{知\zhglue 识\zhglue 卡}
\end{itemize}

\subsection{\zh{五\zhglue 、\zhglue 交\zhglue 付\zhglue 物}}\label{ux4e94ux4ea4ux4ed8ux7269}

{\def\LTcaptype{none} % do not increment counter
\begin{longtable}[]{@{}lll@{}}
\toprule\noalign{}
\zh{交\zhglue 付\zhglue 物} & \zh{位\zhglue 置} & \zh{验\zhglue 收} \\
\midrule\noalign{}
\endhead
\bottomrule\noalign{}
\endlastfoot
ReAct \zh{循\zhglue 环} & \passthrough{\codezh{{\zh{手}}{\zh{写}}/}} & \zh{可\zhglue 跑\zhglue 通\zhglue 多\zhglue 步\zhglue 任\zhglue 务} \\
self-judge \zh{协\zhglue 议} & \passthrough{\codezh{{\zh{手}}{\zh{写}}/}}\zh{（}R10\zh{）} & \zh{有\zhglue 明\zhglue 确\zhglue 的\zhglue 判\zhglue 定\zhglue 时\zhglue 机\zhglue 与\zhglue 阈\zhglue 值\zhglue 设\zhglue 计\zhglue 说\zhglue 明} \\
\zh{轨\zhglue 迹\zhglue 日\zhglue 志} & \passthrough{\codezh{{\zh{证}}{\zh{据}}/trajectories/}} & \zh{每\zhglue 任\zhglue 务\zhglue 一\zhglue 份\zhglue ，\zhglue 含\zhglue 步\zhglue 数\zhglue 与\zhglue 结\zhglue 果} \\
\zh{消\zhglue 融\zhglue 数\zhglue 据} & \passthrough{\codezh{{\zh{证}}{\zh{据}}/}} & \zh{三\zhglue 组\zhglue 成\zhglue 功\zhglue 率} + \zh{误\zhglue 差\zhglue 棒} \\
\zh{报\zhglue 告} & \passthrough{\codezh{{\zh{报}}{\zh{告}}.md}} & \zh{必\zhglue 须\zhglue 回\zhglue 答}''\zh{何\zhglue 时} self-judge \zh{有\zhglue 害}'' \\
\end{longtable}
}

\subsection{\zh{六\zhglue 、\zhglue 讲\zhglue 课\zhglue 锚\zhglue 点}}\label{ux516dux8bb2ux8bfeux951aux70b9}

\begin{itemize}
\tightlist
\item
  \textbf{\zh{讲\zhglue 义}}\zh{：}\passthrough{\codezh{M8-{\zh{智}}{\zh{能}}{\zh{体}}.md}}\zh{（}W13 \zh{展\zhglue 开\zhglue ）}
\item
  \textbf{\zh{对\zhglue 标}}\zh{：}\textbf{\zh{李\zhglue 宏\zhglue 毅\zhglue 《}Context Engineering\zh{》\zhglue 专\zhglue 讲}}\zh{（}Agent \zh{时\zhglue 代\zhglue 的\zhglue 关\zhglue 键\zhglue 技\zhglue 术\zhglue ）}\sym{·} Berkeley LLM Agents MOOC\zh{（}ReAct \zh{原\zhglue 始\zhglue 论\zhglue 文\zhglue 那\zhglue 讲\zhglue ）}\sym{·} HF Agents Course\zh{（\zhglue 三\zhglue 框\zhglue 架\zhglue 对\zhglue 照} + \zh{观\zhglue 测\zhglue 性}/\zh{评\zhglue 测} bonus unit\zh{）}\sym{·} \textbf{rasbt ch5\zh{（}self-refinement\zh{）}} \sym{·} \passthrough{\lstinline!nanochat/execution.py!}+\passthrough{\lstinline!tests/test\_execution.py!}\zh{（}\textbf{\zh{沙\zhglue 箱\zhglue 执\zhglue 行\zhglue 的\zhglue 最\zhglue 小\zhglue 实\zhglue 现}}\zh{，\zhglue 可\zhglue 读\zhglue 接\zhglue 口\zhglue 不\zhglue 可\zhglue 抄\zhglue 实\zhglue 现\zhglue ）}
\item
  \textbf{\zh{搭\zhglue 桥}}\zh{：}M5 \zh{的} KV cache\zh{（\zhglue 长\zhglue 上\zhglue 下\zhglue 文\zhglue 的\zhglue 显\zhglue 存\zhglue 代\zhglue 价\zhglue ）}\sym{→} \zh{为\zhglue 什\zhglue 么\zhglue 长\zhglue 程} Agent \zh{必\zhglue 须\zhglue 做}\textbf{\zh{上\zhglue 下\zhglue 文\zhglue 压\zhglue 缩}}
\end{itemize}

\subsection{\zh{七\zhglue 、\zhglue 代\zhglue 码\zhglue 级\zhglue 提\zhglue 问\zhglue 示\zhglue 例}}\label{ux4e03ux4ee3ux7801ux7ea7ux63d0ux95eeux793aux4f8b}

\begin{enumerate}
\def\labelenumi{\arabic{enumi}.}
\tightlist
\item
  \zh{长\zhglue 程\zhglue 任\zhglue 务\zhglue 失\zhglue 败\zhglue 率\zhglue 随\zhglue 步\zhglue 数\zhglue 指\zhglue 数\zhglue 上\zhglue 升\zhglue 。}\textbf{\zh{分\zhglue 解}}\zh{：\zhglue 是\zhglue 每\zhglue 步\zhglue 错\zhglue 误\zhglue 率\zhglue 问\zhglue 题\zhglue ，\zhglue 还\zhglue 是\zhglue 错\zhglue 误\zhglue 会\zhglue 累\zhglue 积\zhglue 传\zhglue 播\zhglue ？}
\item
  self-judge \zh{什\zhglue 么\zhglue 时\zhglue 候}\textbf{\zh{一\zhglue 定}}\zh{无\zhglue 效\zhglue ？\zhglue （\zhglue 提\zhglue 示\zhglue ：\zhglue 当\zhglue 模\zhglue 型\zhglue 无\zhglue 法\zhglue 区\zhglue 分}''\zh{看\zhglue 起\zhglue 来\zhglue 对}''\zh{和}''\zh{真\zhglue 的\zhglue 对}''\zh{时\zhglue ）}
\item
  \zh{上\zhglue 下\zhglue 文\zhglue 压\zhglue 缩\zhglue （\zhglue 截\zhglue 断}/\zh{摘\zhglue 要\zhglue ）\zhglue 会\zhglue 引\zhglue 入\zhglue 什\zhglue 么\zhglue 新\zhglue 失\zhglue 败\zhglue 模\zhglue 式\zhglue ？}
\item
  \zh{为\zhglue 什\zhglue 么}''\zh{工\zhglue 具\zhglue 返\zhglue 回\zhglue 值\zhglue 要\zhglue 结\zhglue 构\zhglue 化}''\zh{能\zhglue 显\zhglue 著\zhglue 提\zhglue 升\zhglue 成\zhglue 功\zhglue 率\zhglue ？\zhglue （\zhglue 从} token \zh{分\zhglue 布\zhglue 角\zhglue 度\zhglue 解\zhglue 释\zhglue ）}
\item
  \textbf{\zh{【\zhglue 下\zhglue 注\zhglue 】}} \zh{预\zhglue 测} self-judge \zh{在\zhglue 简\zhglue 单\zhglue 任\zhglue 务} vs \zh{困\zhglue 难\zhglue 任\zhglue 务\zhglue 上\zhglue 的\zhglue 效\zhglue 果\zhglue 差\zhglue 异\zhglue ，\zhglue 先\zhglue 写\zhglue 再\zhglue 测\zhglue 。}
\end{enumerate}

\subsection{\zh{八\zhglue 、\zhglue 风\zhglue 险\zhglue 与\zhglue 提\zhglue 醒}}\label{ux516bux98ceux9669ux4e0eux63d0ux9192}

\begin{itemize}
\tightlist
\item
  \emo{🚨} \textbf{\zh{不\zhglue 要\zhglue 用\zhglue 框\zhglue 架\zhglue 起\zhglue 步}}\zh{：\zhglue 手\zhglue 写\zhglue 一\zhglue 遍} ReAct \zh{再\zhglue 去\zhglue 看} LangGraph/smolagents\zh{，\zhglue 否\zhglue 则\zhglue 学\zhglue 不\zhglue 到} harness \zh{的\zhglue 本\zhglue 质\zhglue 。}
\item
  \sym{⚠️} self-judge \zh{会\zhglue 显\zhglue 著\zhglue 增\zhglue 加} token \zh{消\zhglue 耗\zhglue 与\zhglue 时\zhglue 间} \sym{→} \zh{算\zhglue 力\zhglue 预\zhglue 算\zhglue 要\zhglue 写\zhglue 进\zhglue 实\zhglue 验\zhglue 卡\zhglue 。}
\item
  \sym{⚠️} \zh{轨\zhglue 迹\zhglue 日\zhglue 志\zhglue 可\zhglue 能\zhglue 很\zhglue 大} \sym{→} \zh{只\zhglue 保\zhglue 留\zhglue 结\zhglue 构\zhglue 化\zhglue 摘\zhglue 要\zhglue （}thought \zh{截\zhglue 断} + \zh{结\zhglue 果\zhglue ）\zhglue ，\zhglue 避\zhglue 免\zhglue 污\zhglue 染\zhglue 仓\zhglue 库\zhglue 体\zhglue 积\zhglue 。}
\end{itemize}

\end{document}